\BourbakiExerciseGroup

\begin{enumerate}
\def\labelenumi{\arabic{enumi})}
\tightlist
\item
  Let \(u(\mathbf{X}) = \sum_{n=0}^{\infty} \alpha_n \mathbf{X}^n\) be a formal power series over a commutative field \(\mathbf{K}\) .
\end{enumerate}

\begin{enumerate}
\def\labelenumi{\alph{enumi})}
\tightlist
\item
  For u to be a rational fraction of \(\mathbf{K}(\mathbf{X})\) it is necessary and sufficient that there should exist a finite sequence of elements \((\lambda_{i})_{1\leq i\leq q}\) of elements of K, not all zero, and an integer \(d\geq0\) such that for all \(n\geq d\) we have
\end{enumerate}

\[
\lambda_ {1} \alpha_ {n} + \lambda_ {2} \alpha_ {n + 1} + \dots + \lambda_ {q} \alpha_ {n + q - 1} = 0.
\]

\begin{enumerate}
\def\labelenumi{\alph{enumi})}
\setcounter{enumi}{1}
\tightlist
\item
  Put
\end{enumerate}

\[
\mathrm{H} _ {n} ^ {(k)} = \left| \begin{array}{c c c c} \alpha_ {n} & \alpha_ {n + 1} & \dots & \alpha_ {n + k - 1} \\ \alpha_ {n + 1} & \alpha_ {n + 2} & \dots & \alpha_ {n + k} \\ \alpha_ {n + 2} & \alpha_ {n + 3} & \dots & \alpha_ {n + k + 1} \\ \vdots & \vdots & \ddots & \vdots \\ \alpha_ {n + k - 1} & \alpha_ {n + k} & \dots & \alpha_ {n + 2 k - 2} \end{array} \right|
\]

(« Hankel determinant ») show that if \(\mathrm{H}_{d+j}^{(q+1)} = 0\) and \(\mathrm{H}_{d+j}^{(q)} \neq 0\) for every integer \(j \geqslant 0\) , then \(u(\mathbf{X})\) is a rational fraction (use \(a\) )).

\begin{enumerate}
\def\labelenumi{\alph{enumi})}
\setcounter{enumi}{2}
\tightlist
\item
  Prove the identity
\end{enumerate}

\[
\mathrm{H} _ {n} ^ {(k)} \mathrm{H} _ {n + 2} ^ {(k)} - \mathrm{H} _ {n} ^ {(k + 1)} \mathrm{H} _ {n + 2} ^ {(k - 1)} = (\mathrm{H} _ {n + 1} ^ {(k)}) ^ {2}
\]

(cf.~III, p.~639, Exercise 10). Deduce that if \(\mathrm{H}_{m+j}^{(k+1)} = 0\) for \(0 \leqslant j \leqslant r-1\) , then the \(r\) determinants \(\mathrm{H}_{m+j}^{(k)}\) where \(1 \leqslant j \leqslant r\) are either all zero or all non-zero.

\(d\) ) Deduce from \(b\) ) and \(c\) ) that for \(u(\mathbf{X})\) to be a rational fraction it is necessary and sufficient that there exist two integers \(d\) and \(q\) such that \(\mathrm{H}_{d + j}^{(q + 1)} = 0\) for all integers \(j\geqslant 0\)

\begin{enumerate}
\def\labelenumi{\alph{enumi})}
\setcounter{enumi}{4}
\tightlist
\item
  Let \(k\) be a subfield of \(\mathbf{K}\) , and identify \(\mathbf{K}(\mathbf{X})\) , \(k(\mathbf{X})\) and \(k[[\mathbf{X}]]\) with subrings of \(\mathbf{K}((\mathbf{X}))\) . Prove that \(k[[\mathbf{X}]] \cap \mathbf{K}(\mathbf{X}) = k[[\mathbf{X}]] \cap k(\mathbf{X})\) .
\end{enumerate}

\begin{enumerate}
\def\labelenumi{\arabic{enumi})}
\setcounter{enumi}{1}
\tightlist
\item
  Let \(a_1, a_2, \ldots, a_p\) be integers \(> 0\) ; denote by \(\alpha_n\) the number of finite sequences \((x_i)_{1 \leqslant i \leqslant p}\) of integers \(\geqslant 0\) satisfying the equation
\end{enumerate}

\[
a _ {1} x _ {1} + a _ {2} x _ {2} + \dots + a _ {p} x _ {p} = n.
\]

Show that the formal power series \(\sum_{n=0}^{\infty} \alpha_n X^n\) (over \(\mathbf{Q}\) ) is the expansion of the rational fraction

\[
\frac {1}{(1 - \mathrm{X} ^ {a _ {1}}) (1 - \mathrm{X} ^ {a _ {2}}) \dots (1 - \mathrm{X} ^ {a _ {p}})}.
\]

\begin{enumerate}
\def\labelenumi{\arabic{enumi})}
\setcounter{enumi}{2}
\tightlist
\item
  Let \(\mathbf{F}\) be a finite set of integers \(> 0\) , and let \(\beta_{n}\) be the number of finite sequences \((x_{i})\) of at most \(n\) terms, all of whose terms lie in \(\mathbf{F}\) and such that \(\sum_{i} x_{i} = n\) . Show that the formal power series \(\sum_{n=0}^{\infty} \beta_n X^n\) over \(Q\) is the expansion of the rational fraction
\end{enumerate}

\[
\frac {1}{1 - \sum_ {a \in F} X ^ {a}}.
\]

\begin{enumerate}
\def\labelenumi{\arabic{enumi})}
\setcounter{enumi}{3}
\item
  Let E be a vector space with an infinite basis over a field K of characteristic 2; let B be the exterior algebra \(\Lambda E\) on this space, which is a commutative ring. Give an example of a formal power series \(u \in B[[X]]\) such that \(u^2 = 0\) but such that there exists no element \(\gamma \neq 0\) of B for which \(\gamma u = 0\) .
\item
  Let \(\mathbf{E} = \mathbf{A}[\mathbf{X}_1, \dots, \mathbf{X}_p]\) , \(\mathbf{F} = \mathbf{A}[[\mathbf{X}_1, \dots, \mathbf{X}_p]]\) .
\end{enumerate}

\begin{enumerate}
\def\labelenumi{\alph{enumi})}
\item
  If D is a Z-derivation of E into F, we have \(\omega(Du) \geqslant \omega(u) - 1\) for every non-zero \(u\) in E. b) If D' is a Z-derivation of F, then \(\omega(D'v) \geqslant \omega(v') - 1\) for every non-zero \(v\) in F.
\item
  If E is equipped with the topology induced by that of F, then D and D' are continuous. d) Every Z-derivation of E into F extends in a unique fashion to a Z-derivation of F. e) Let D be the F-module of A-derivations of F. If D ∈ D and DX\_i = u\_i, then D = ∑\_\{i=1\}\^{}\{p\} u\_i D\_i.
\end{enumerate}

\(f)\) \((\mathbf{D}_1,\dots ,\mathbf{D}_p)\) is a basis of the F-module \(\mathcal{D}\) .

\begin{enumerate}
\def\labelenumi{\arabic{enumi})}
\setcounter{enumi}{5}
\tightlist
\item
  Let \((u_{\lambda})_{\lambda \in L}\) be a family of elements of \(\mathbf{A}[[\mathbf{X}_i]]_{i\in I}\) . If the family \((1 + u_{\lambda})_{\lambda \in L}\) is multipliable, the family \((u_{\lambda})_{\lambda \in L}\) is summable.
\end{enumerate}

* 7) If A is reduced, then so is \(\mathbf{A}[[(\mathbf{X}_i)_{i\in I}]]\) . *

\begin{enumerate}
\def\labelenumi{\arabic{enumi})}
\setcounter{enumi}{7}
\tightlist
\item
  We denote by \((1 + \mathbf{X})^{\mathrm{T}}\) the element \(\exp (\mathrm{T}\log (1 + \mathbf{X}))\) of \(\mathbf{Q}[[\mathbf{X},\mathrm{T}]]\) .
\end{enumerate}

\begin{enumerate}
\def\labelenumi{\alph{enumi})}
\tightlist
\item
  Show that \((1 + \mathbf{X})^{\mathrm{T}} = \sum_{n\geqslant 0}\frac{\mathrm{T}(\mathrm{T} - 1)\dots(\mathrm{T} - n + 1)}{n!}\mathbf{X}^{n}\) . (The coefficient of \(\mathbf{X}^n\) is a
\end{enumerate}

polynomial in T whose value is known for \(\mathbf{T} = n\in \mathbf{N}\) .)

\begin{enumerate}
\def\labelenumi{\alph{enumi})}
\setcounter{enumi}{1}
\tightlist
\item
  Show that \((1 + \mathbf{X})^{\mathrm{T}}(1 + \mathbf{X})^{\mathrm{T}} = (1 + \mathbf{X})^{\mathrm{T} + \mathrm{T}}\) . Write down explicitly the identity between binomial coefficients so obtained. Show that \((1 + \mathbf{X} + \mathbf{Y} + \mathbf{XY})^{\mathrm{T}} = (1 + \mathbf{X})^{\mathrm{T}}(1 + \mathbf{Y})^{\mathrm{T}}\) .
\end{enumerate}

